{-
Introduction to Type Theory (1) - Verona

using agda

What is the Type Theory?
- a programming language
  with a very strong typrev {ℕ} (1 ∷ 2 ∷ 3 ∷ [])e system (dependent types)
- a (constructive) logic, intuitionistic
- alternative foundation of Mathematics 
  Alternative to axiomatic set theory (ZFC)

Features :

- typing is a judgement
3 ∈ ℕ   ,  3 : ℕ
TT statically typed, typing is judgement
univalence axiom
∪ ∩ ⊆ not on types, but on subsetes of a given type
- Logic is emergent (propositions as types)
- Functions are primitive 
  R : A → B → Prop


Plan :
Monday : intro + dependent types in programming
Tuesday : logic , PaT
Wednesday : Inductive relations, meta logic 
Friday : Cubical Type Theory (Homotopy Type Theory)
-}

-- pattern match on m, C-c C-c
-- check shed C-c C-SPC

open import Data.Nat
open import Data.List

-- reverse
-- rev (1 ∷ 2 ∷ 3 ∷ []) = 3 ∷ 2 ∷ 1 ∷ [] -- \::

-- List : Set → Set 
-- polymorphism is a special case of dependent types

snoc : {A : Set} → List A → A → List A
snoc [] a = a ∷ []
snoc (a ∷ as) b = a ∷ snoc as b

rev : {A : Set} → List A → List A
rev [] = []
rev (a ∷ as) = snoc (rev as) a

length' : {A : Set} → List A → ℕ
length' [] = zero
length' (a ∷ as) = suc (length as)

variable A B C : Set

open import Data.Maybe

_!!_ : List A → ℕ → Maybe A
[] !! n = nothing
(a ∷ as) !! zero = just a
(a ∷ as) !! suc n = as !! n

open import Data.Vec
-- Vec : Set → ℕ → Set
-- a proper dependent type (not just polymorphic)

{-
data Fin : ℕ → Set where
  zero : {n : ℕ} → Fin (suc n)
  suc  : {n : ℕ} (i : Fin n) → Fin (suc n)
-}
-- Fin 0 = []
-- Fin 1 = [ zero ]
-- Fin 2 = [ zero , suc zero ]
-- Fin 3 = [ zero , suc zero , suc (suc zero) ]
-- Fin n = { 0 , 1 , .., n-1 }

-- Fin : ℕ → Set

snoc-v : {n : ℕ} → Vec A n → A → Vec A (suc n)
snoc-v [] a = a ∷ []
snoc-v (a ∷ as) b = a ∷ snoc-v as b

rev-v : {n : ℕ} → Vec A n → Vec A n
rev-v [] = []
rev-v (a ∷ as) = snoc-v (rev-v as) a

open import Data.Fin hiding (_+_)

_!!-v_ : {n : ℕ} → Vec A n → Fin n → A
(a ∷ as) !!-v zero = a
(a ∷ as) !!-v suc i = as !!-v i

map-v : {n : ℕ} → (A → B) → (Vec A n) → (Vec B n)
map-v f [] = []
map-v f (a ∷ as) = f a ∷ map-v f as

enum : (n : ℕ) → Vec (Fin n) n
enum zero = []
enum (suc n) = zero ∷ (map-v suc (enum n)) -- C-c C-. show GOAL/HAVE

max : {n : ℕ} → Fin (suc n)
max {zero} = zero
max {suc n} = suc (max {n}) -- C-c C-r

emb : {n : ℕ} → Fin n → Fin (suc n)
emb zero = zero
emb (suc i) = suc (emb i)

inv : {n : ℕ} → Fin n → Fin n
inv zero = max
inv (suc i) = emb (inv i) -- left as an exercise

--
-- vectors and matrices 

Vector : ℕ → Set {- Vec n is an n-dimensional vector -}
Vector m = Vec ℕ m

Matrix : ℕ → ℕ → Set {- Matrix m n is an m x n Matrix -}
Matrix m n = Vec (Vector n) m

{- multiplication with a scalar -}
_*v_ : {n : ℕ} → ℕ → Vector n → Vector n
k *v [] = []
k *v (x ∷ xs) = k * x ∷ (k *v xs)

v1 : Vector 3
v1 = 1 ∷ 2 ∷ 3 ∷ []

test1 : Vector 3
test1 = 2 *v v1
{- 2 ∷ 4 ∷ 6 ∷ [] -}

{- addition of vectors -}
_+v_ : {n : ℕ} → Vector n → Vector n → Vector n
[] +v [] = []
(x ∷ xs) +v (y ∷ ys) = x + y ∷  (xs +v ys)

v2 : Vector 3
v2 = 2 ∷ 3 ∷ 0 ∷ []

test2 : Vector 3
test2 = v1 +v v2
{- 3 ∷ 5 ∷ 3 ∷ [] -}

zeros : {n : ℕ} → Vector n
zeros {zero} = []
zeros {suc n} = zero ∷ zeros {n}

{- multiplication of a vector and a matrix -}
_*vm_ : {m n : ℕ} → Vector m → Matrix m n → Vector n
[] *vm [] = zeros
(x ∷ xs) *vm (ys ∷ xys) = (x *v ys) +v (xs *vm xys)

id3 : Matrix 3 3
id3 = (1 ∷ 0 ∷ 0 ∷ []) 
    ∷ (0 ∷ 1 ∷ 0 ∷ []) 
    ∷ (0 ∷ 0 ∷ 1 ∷ []) 
    ∷ []

m3 : Matrix 3 3
m3 =  (1 ∷ 2 ∷ 3 ∷ []) 
    ∷ (4 ∷ 5 ∷ 6 ∷ []) 
    ∷ (7 ∷ 8 ∷ 9 ∷ []) 
    ∷ []

test3 : Vector 3
test3 = v1 *vm m3
{- 30 ∷ 36 ∷ 42 ∷ [] -}

{- matrix multiplication -}
-- _*mm_ : {l m n : ℕ} → Matrix l m → Matrix m n → Matrix l n
-- xys *mm xys' = ?

inv3 : Matrix 3 3
inv3 = (0 ∷ 0 ∷ 1 ∷ []) 
     ∷ (0 ∷ 1 ∷ 0 ∷ []) 
     ∷ (1 ∷ 0 ∷ 0 ∷ []) 
     ∷ []

-- test4 : Matrix 3 3
-- test4 = inv3 *mm inv3
{-
(30 ∷ 36 ∷ 42 ∷ []) ∷
(66 ∷ 81 ∷ 96 ∷ []) ∷
(102 ∷ 126 ∷ 150 ∷ [])
∷ []
-}

{- implement matrix transposition, i.e. swap rows and columns. -}

-- transp : {m n : ℕ} → Matrix m n → Matrix n m
-- transp xys = ? 

rect : Matrix 2 3
rect = (1 ∷ 2 ∷ 3 ∷ [])
     ∷ (4 ∷ 5 ∷ 6 ∷ [])
     ∷ []

-- test5 : Matrix 3 2
-- test5 = transp rect
{-
(1 ∷ 4 ∷ []) ∷
(2 ∷ 5 ∷ []) ∷
(3 ∷ 6 ∷ []) ∷ []
-}


